\section{Giới thiệu và mục tiêu}

Bài báo cáo này trình bày trực quan hóa hiệu năng của ba thuật toán sắp xếp điển hình gồm \textbf{Insertion Sort}, \textbf{Quick Sort} và \textbf{Counting Sort} trên bốn kiểu dữ liệu đầu vào:
\begin{itemize}[leftmargin=1.2cm]
    \item Random.
    \item Reversed.
    \item Nearly Sorted.
    \item Many Duplicates.
\end{itemize}

Mục tiêu chính của bài là:
\begin{itemize}[leftmargin=1.2cm]
    \item Xây dựng các biểu đồ đúng chuẩn trực quan hóa dữ liệu để so sánh thuật toán theo kích thước đầu vào $n$.
    \item So sánh và phân tích khác biệt khi dùng \textbf{thang tuyến tính} và \textbf{thang log trên trục $y$}.
    \item Rút ra insight về hành vi thuật toán theo từng đặc trưng dữ liệu.
    \item Đánh giá các kỹ thuật trực quan được sử dụng và lý do lựa chọn.
\end{itemize}

\subsection{Phạm vi và giả định}

Bài làm sử dụng trực tiếp dữ liệu đo hiệu năng đã được cung cấp trong thư mục \texttt{results}.
Không chạy lại benchmark và không phát sinh dữ liệu mới ngoài bộ dữ liệu đã có.

\subsection{Câu hỏi phân tích}

Các biểu đồ trong báo cáo trả lời các câu hỏi sau:
\begin{itemize}[leftmargin=1.2cm]
    \item Thuật toán nào nhanh hơn theo từng mốc $n$ và từng kiểu dữ liệu?
    \item Sự khác biệt giữa các thuật toán thay đổi như thế nào khi chuyển từ linear sang log-$y$?
    \item Tại sao một số thuật toán có lợi thế rõ rệt trên dữ liệu đặc thù (nearly sorted, many duplicates)?
\end{itemize}
